% !TEX spellcheck = en_US



% ~~~~~~~~~~~~~~~~~~~~~~~~~~~~~~~~~~~~~~~ V
% (NC) 2026 Haluk Bingol
% github.com/halukbingol/LaTeX-Templates
%
% Licensees may copy, distribute, display, and perform the work and make derivative works 
% and remixes based on it only for non-commercial purposes. 
% ~~~~~~~~~~~~~~~~~~~~~~~~~~~~~~~~~~~~~~~ A




% =====================================================================
%  Internationalization in Java -- Sophomore Level: A Localized Application
%  ---------------------------------------------------------------------
%  Built on Template/hblibTemplate.tex with the libraries in ../../hblib/.
%  Programs, terminal sessions and their real outputs are in codeoutput/:
%     codeoutput/<Name>.java  a program        codeoutput/<Name>.sh   a session
%     codeoutput/<Name>/      the session's project folder
%     codeoutput/<Name>.txt   what it printed, made by:  bash run_examples.sh
% =====================================================================

\documentclass[aspectratio=169,10pt,compress]{beamer}

% ~~~~~~~~~~~~~~~~~~~~~~~~~~~~~~~~~~~~~~~~~~~~~~~~~~~~~~~~~~~~~~~~~ V hblib
	\usepackage{../../hblib/hblibPresentation} % lib presentation: theme, i/N, \hSection, algorithm2e, ...
	\usepackage{../../hblib/hblibCodeoutput} % lib code + output: \lstcode, \lstterminal, ...
	\usepackage{../../hblib/hblibDatastructures} % lib data structures: \hString, \hStack, \hQueue
% ~~~~~~~~~~~~~~~~~~~~~~~~~~~~~~~~~~~~~~~~~~~~~~~~~~~~~~~~~~~~~~~~~ A hblib


% xcolor, array (and the L{width} column), listings, tikz and algorithm2e come from the hblib libraries
\usepackage[T1]{fontenc}
\usepackage{lmodern}
\usepackage{booktabs}

\definecolor{gitred}{RGB}{240,80,50}
\newcommand{\win}{\textbf{Windows}}
\newcommand{\mac}{\textbf{macOS}}

% algorithm2e set-up (the package is loaded by hblibPresentation), as in the template
\SetAlCapFnt{\small}
\SetAlCapNameFnt{\small}
\DontPrintSemicolon

% .properties files (resource bundles)
\lstdefinelanguage{properties}{morecomment=[l]{\#}, morecomment=[l]{!}}

% ---------------------------------------------------------------------
\title{Internationalization in Java}
\subtitle{Sophomore Level: A Localized Application}
\author[Bingol]{Haluk O. Bingol}

\institute[FCIS]{
    Faculty of Computer and Information Sciences\\
    Yeditepe University
}
\date{
	{\tiny \hbTimeStamp}
}

\begin{document}




% ~~~~~~~~~~~~~~~~~~~~~~~~~~~~~~~~~~~~~~ V frm
\begin{frame}[t,fragile]
  \titlepage
\end{frame}
% ~~~~~~~~~~~~~~~~~~~~~~~~~~~~~~~~~~~~~~ A frm




% ~~~~~~~~~~~~~~~~~~~~~~~~~~~~~~~~~~~~~~ V frm
\begin{frame}[t,fragile] \frametitle{Outline}
  \begin{columns}[T]
    \begin{column}{0.48\textwidth}
      \tableofcontents[sections={1-3}, hideallsubsections]
    \end{column}
    \begin{column}{0.48\textwidth}
      \tableofcontents[sections={4-20}, hideallsubsections]
    \end{column}
  \end{columns}
\end{frame}
% ~~~~~~~~~~~~~~~~~~~~~~~~~~~~~~~~~~~~~~ A frm

% ~~~~~~~~~~~~~~~~~~~~~~~~~~~~~~~~~~~~~~ V
\hSection[Project]{A Localized Application}
% ~~~~~~~~~~~~~~~~~~~~~~~~~~~~~~~~~~~~~~ A

% ~~~~~~~~~~~~~~~~~~~~~~~~~~~~~~~~~~~~~~ V
\subsection[Layout]{Project layout}
% ~~~~~~~~~~~~~~~~~~~~~~~~~~~~~~~~~~~~~~ A




% ~~~~~~~~~~~~~~~~~~~~~~~~~~~~~~~~~~~~~~ V frm
\begin{frame}[t,fragile] \frametitle{The project}
  \lstterminal[basicstyle=\ttfamily\scriptsize, lastline=6]{codeoutput/Build.txt}
  \begin{itemize}\small
    \item Bundles are \textbf{resources}: they live in a package (\texttt{edu.yeditepe.shop.i18n}) next to the code.
    \item The \textbf{base name} is the full name: \texttt{edu.yeditepe.shop.i18n.Messages}.
    \item One file per language; the base file (no suffix) is English and has every key.
  \end{itemize}
\end{frame}
% ~~~~~~~~~~~~~~~~~~~~~~~~~~~~~~~~~~~~~~ A frm




% ~~~~~~~~~~~~~~~~~~~~~~~~~~~~~~~~~~~~~~ V frm
\begin{frame}[t,fragile] \frametitle{The bundles}
  \lstcode[basicstyle=\ttfamily\tiny, numbers=none]{codeoutput/Shop/src/edu/yeditepe/shop/i18n/Messages.properties}{properties}
  \lstcode[basicstyle=\ttfamily\tiny, numbers=none]{codeoutput/Shop/src/edu/yeditepe/shop/i18n/Messages_tr.properties}{properties}
  \lstcode[basicstyle=\ttfamily\tiny, numbers=none]{codeoutput/Shop/src/edu/yeditepe/shop/i18n/Messages_de.properties}{properties}
\end{frame}
% ~~~~~~~~~~~~~~~~~~~~~~~~~~~~~~~~~~~~~~ A frm




% ~~~~~~~~~~~~~~~~~~~~~~~~~~~~~~~~~~~~~~ V frm
\begin{frame}[t,fragile] \frametitle{The code: no text, no format in it}
  \lstcode[basicstyle=\ttfamily\tiny, firstline=13, lastline=35, firstnumber=13]{codeoutput/Shop/src/edu/yeditepe/shop/Main.java}{Java}
  \vspace{0.1em}
  \small Every text comes from the bundle; every number, amount and date is formatted for the locale \texttt{l}.
  The Turkish \texttt{items} has no \texttt{1\#} case: \emph{1 ürün}, \emph{3 ürün}.
\end{frame}
% ~~~~~~~~~~~~~~~~~~~~~~~~~~~~~~~~~~~~~~ A frm

% ~~~~~~~~~~~~~~~~~~~~~~~~~~~~~~~~~~~~~~ V
\subsection[Build]{Build and run}
% ~~~~~~~~~~~~~~~~~~~~~~~~~~~~~~~~~~~~~~ A




% ~~~~~~~~~~~~~~~~~~~~~~~~~~~~~~~~~~~~~~ V frm
\begin{frame}[t,fragile] \frametitle{Building: the bundles must be copied}
  \lstterminal[basicstyle=\ttfamily\tiny]{codeoutput/Build.txt}

  \vspace{0.1em}
  \small
  \texttt{javac} only writes \texttt{.class} files. The \texttt{.properties} files must be copied into \texttt{out}
  (Maven and Gradle do this for \texttt{src/main/resources}). \win: \texttt{mkdir} and \texttt{copy} with \texttt{\textbackslash}.
\end{frame}
% ~~~~~~~~~~~~~~~~~~~~~~~~~~~~~~~~~~~~~~ A frm




% ~~~~~~~~~~~~~~~~~~~~~~~~~~~~~~~~~~~~~~ V frm
\begin{frame}[t,fragile] \frametitle{In a JAR}
  \lstterminal[basicstyle=\ttfamily\tiny]{codeoutput/Jar.txt}

  \vspace{0.1em}
  \small
  The bundles travel inside the JAR. Without an argument the program uses the default locale:
  \texttt{de-AT} has no file, so it uses \texttt{Messages\_de}; the money format is Austrian (\texttt{TRY 3.749,70}).
\end{frame}
% ~~~~~~~~~~~~~~~~~~~~~~~~~~~~~~~~~~~~~~ A frm

% ~~~~~~~~~~~~~~~~~~~~~~~~~~~~~~~~~~~~~~ V
\hSection[Lookup]{How Bundles Are Found}
% ~~~~~~~~~~~~~~~~~~~~~~~~~~~~~~~~~~~~~~ A




% ~~~~~~~~~~~~~~~~~~~~~~~~~~~~~~~~~~~~~~ V frm
\begin{frame}[t,fragile] \frametitle{The search order}
  \begin{columns}[T]
    \begin{column}{0.48\textwidth}
      \small Request \texttt{de-CH}, default locale \texttt{tr-TR}:
      \begin{enumerate}\small
        \item \texttt{Messages\_de\_CH}
        \item \texttt{Messages\_de}
        \item \texttt{Messages\_tr\_TR} \hRemark{(default locale!)}
        \item \texttt{Messages\_tr}
        \item \texttt{Messages}
      \end{enumerate}
      \small The first file that \textbf{exists} is returned. Its \emph{parents} are the later files of the
      same language and the base file, so missing keys are still found.
    \end{column}
    \begin{column}{0.48\textwidth}
      \begin{block}{Only when nothing matches}
        \small Steps 3--4 are tried only if neither \texttt{\_de\_CH} nor \texttt{\_de} exists.
        Then a French user on a Turkish server gets \textbf{Turkish}.
      \end{block}
      \begin{block}{Each file can be}
        \small a \texttt{.properties} file, or a class \texttt{Messages\_de extends ListResourceBundle}
        (for values that are not text, e.g.\ images or objects). Classes are tried first.
      \end{block}
    \end{column}
  \end{columns}
\end{frame}
% ~~~~~~~~~~~~~~~~~~~~~~~~~~~~~~~~~~~~~~ A frm




% ~~~~~~~~~~~~~~~~~~~~~~~~~~~~~~~~~~~~~~ V frm
\begin{frame}[t,fragile] \frametitle{The default-locale trap}
  \lstterminal[basicstyle=\ttfamily\scriptsize]{codeoutput/Trap.txt}

  \vspace{0.1em}
  \small
  The same request \texttt{fr-FR}: on an English server the texts are English (base file); on a Turkish server
  they are \textbf{Turkish}, while money and date are French. Mixed languages on one screen.
\end{frame}
% ~~~~~~~~~~~~~~~~~~~~~~~~~~~~~~~~~~~~~~ A frm




% ~~~~~~~~~~~~~~~~~~~~~~~~~~~~~~~~~~~~~~ V frm
\begin{frame}[t,fragile] \frametitle{Turning the fallback off}
  \begin{columns}[T]
    \begin{column}{0.58\textwidth}
      \lstcode[basicstyle=\ttfamily\tiny]{codeoutput/Lookup/Lookup.java}{Java}
    \end{column}
    \begin{column}{0.40\textwidth}
      \uncover<2->{\lstterminal[basicstyle=\ttfamily\tiny]{codeoutput/Lookup.txt}}
    \end{column}
  \end{columns}

  \vspace{0.1em}
  \small \texttt{Control.getNoFallbackControl} skips the default locale: unknown languages get the base file,
  whatever the server's locale is. Servers that serve many users should use it.
\end{frame}
% ~~~~~~~~~~~~~~~~~~~~~~~~~~~~~~~~~~~~~~ A frm

% ~~~~~~~~~~~~~~~~~~~~~~~~~~~~~~~~~~~~~~ V
\hSection[Files]{Bundle Files and Encodings}
% ~~~~~~~~~~~~~~~~~~~~~~~~~~~~~~~~~~~~~~ A




% ~~~~~~~~~~~~~~~~~~~~~~~~~~~~~~~~~~~~~~ V frm
\begin{frame}[t,fragile] \frametitle{Encoding of \texttt{.properties} files}
  \lstterminal[basicstyle=\ttfamily\tiny]{codeoutput/Encoding.txt}
  \vspace{0.1em}
  \begin{itemize}\small
    \item Since Java 9 bundles are read as \textbf{UTF-8}. If a file is not valid UTF-8, Java reads it as
      ISO-8859-1: a Turkish ISO-8859-9 file shows \texttt{Çýkýþ}.
    \item \texttt{\textbackslash u0131} escapes work everywhere (the only way before Java 9).
    \item Save bundle files as UTF-8 in the editor (\win: Notepad \emph{Save as, Encoding: UTF-8}).
  \end{itemize}
\end{frame}
% ~~~~~~~~~~~~~~~~~~~~~~~~~~~~~~~~~~~~~~ A frm

% ~~~~~~~~~~~~~~~~~~~~~~~~~~~~~~~~~~~~~~ V
\hSection[Messages]{MessageFormat in Depth}
% ~~~~~~~~~~~~~~~~~~~~~~~~~~~~~~~~~~~~~~ A




% ~~~~~~~~~~~~~~~~~~~~~~~~~~~~~~~~~~~~~~ V frm
\begin{frame}[t,fragile] \frametitle{Pitfalls of \texttt{MessageFormat}}
  \begin{columns}[T]
    \begin{column}{0.62\textwidth}
      \lstcode[basicstyle=\ttfamily\tiny]{codeoutput/Pitfalls.java}{Java}
    \end{column}
    \begin{column}{0.36\textwidth}
      \uncover<2->{\lstterminal[basicstyle=\ttfamily\tiny]{codeoutput/Pitfalls.txt}}
    \end{column}
  \end{columns}

  \vspace{0.1em}
  \small A number argument is formatted with grouping: \texttt{2,026}. A single \texttt{'} quotes text, so
  \texttt{It's} loses its apostrophe and \texttt{'\{0\}'} is printed literally. Translators must know these rules.
\end{frame}
% ~~~~~~~~~~~~~~~~~~~~~~~~~~~~~~~~~~~~~~ A frm




% ~~~~~~~~~~~~~~~~~~~~~~~~~~~~~~~~~~~~~~ V frm
\begin{frame}[t,fragile] \frametitle{Argument types}
  \begin{center}\small
  \begin{tabular}{ll}
    \toprule
    \textbf{pattern} & \textbf{meaning} \\
    \midrule
    \texttt{\{0\}}                        & \texttt{toString()}, numbers and dates localized \\
    \texttt{\{0,number\}}, \texttt{\{0,number,integer\}}, \texttt{\{0,number,percent\}} & numbers \\
    \texttt{\{0,number,\#.\#\#\}}          & a \texttt{DecimalFormat} pattern \\
    \texttt{\{0,date,long\}}, \texttt{\{0,time,short\}} & \texttt{java.util.Date} only, not \texttt{java.time} \\
    \texttt{\{0,choice,0\#none|1\#one|1<\{0\} many\}} & select a text by a number \\
    \texttt{''}                           & one apostrophe \\
    \texttt{'\{'}                         & a literal brace \\
    \bottomrule
  \end{tabular}
  \end{center}
  \small Format \texttt{java.time} values with \texttt{DateTimeFormatter} first and pass the \texttt{String}, as \texttt{Main} does.
\end{frame}
% ~~~~~~~~~~~~~~~~~~~~~~~~~~~~~~~~~~~~~~ A frm

% ~~~~~~~~~~~~~~~~~~~~~~~~~~~~~~~~~~~~~~ V
\hSection[Input]{Reading Localized Input}
% ~~~~~~~~~~~~~~~~~~~~~~~~~~~~~~~~~~~~~~ A




% ~~~~~~~~~~~~~~~~~~~~~~~~~~~~~~~~~~~~~~ V frm
\begin{frame}[t,fragile] \frametitle{Parsing numbers and dates safely}
  \begin{columns}[T]
    \begin{column}{0.66\textwidth}
      \lstcode[basicstyle=\ttfamily\tiny, firstline=10, firstnumber=10]{codeoutput/ParseInput.java}{Java}
    \end{column}
    \begin{column}{0.32\textwidth}
      \uncover<2->{\lstterminal[basicstyle=\ttfamily\tiny]{codeoutput/ParseInput.txt}}
    \end{column}
  \end{columns}

  \vspace{0.1em}
  \small \texttt{parse(s, pos)} tells how far it got: \texttt{12,5 kg} is not a number, although plain \texttt{parse}
  would return \texttt{12.5}. Dates in a short style are ambiguous: \texttt{10/15/26} only makes sense in the US.
\end{frame}
% ~~~~~~~~~~~~~~~~~~~~~~~~~~~~~~~~~~~~~~ A frm

% ~~~~~~~~~~~~~~~~~~~~~~~~~~~~~~~~~~~~~~ V
\hSection[Sorting]{Comparing Text}
% ~~~~~~~~~~~~~~~~~~~~~~~~~~~~~~~~~~~~~~ A




% ~~~~~~~~~~~~~~~~~~~~~~~~~~~~~~~~~~~~~~ V frm
\begin{frame}[t,fragile] \frametitle{Collator strength}
  \begin{columns}[T]
    \begin{column}{0.46\textwidth}
      \lstcode[basicstyle=\ttfamily\tiny]{codeoutput/Strength.java}{Java}
    \end{column}
    \begin{column}{0.52\textwidth}
      \uncover<2->{\lstterminal[basicstyle=\ttfamily\tiny]{codeoutput/Strength.txt}}
    \end{column}
  \end{columns}

  \vspace{0.1em}
  \small PRIMARY: only base letters (\texttt{c} = \texttt{ç} is false in Turkish: different letters!).
  SECONDARY: accents count. TERTIARY: case counts too. Use PRIMARY for search, TERTIARY for sorting.
\end{frame}
% ~~~~~~~~~~~~~~~~~~~~~~~~~~~~~~~~~~~~~~ A frm

% ~~~~~~~~~~~~~~~~~~~~~~~~~~~~~~~~~~~~~~ V
\hSection[Testing]{Testing in Many Locales}
% ~~~~~~~~~~~~~~~~~~~~~~~~~~~~~~~~~~~~~~ A




% ~~~~~~~~~~~~~~~~~~~~~~~~~~~~~~~~~~~~~~ V frm
\begin{frame}[t,fragile] \frametitle{One loop, seven locales}
  \lstterminal[basicstyle=\ttfamily\tiny]{codeoutput/AllLocales.txt}

  \vspace{0.1em}
  \small
  Run every test in several locales: at least English, Turkish (case and comma) and German or French (long words, spaces).
  \texttt{de-CH} and \texttt{pt-BR} have no own file: they reuse \texttt{\_de} and the base file.
\end{frame}
% ~~~~~~~~~~~~~~~~~~~~~~~~~~~~~~~~~~~~~~ A frm




% ~~~~~~~~~~~~~~~~~~~~~~~~~~~~~~~~~~~~~~ V frm
\begin{frame}[t,fragile] \frametitle{Pseudo-localization: a small tool}
  \lstcode[basicstyle=\ttfamily\tiny, firstline=8, lastline=22, firstnumber=8]{codeoutput/Pseudo/Pseudo.java}{Java}
  \vspace{0.1em}
  \small A fake language that everyone can still read. Lines 23--30 (not shown) apply \texttt{pseudo} to every
  value of \texttt{Messages.properties} and write \texttt{Messages\_en\_XA.properties}.
\end{frame}
% ~~~~~~~~~~~~~~~~~~~~~~~~~~~~~~~~~~~~~~ A frm




% ~~~~~~~~~~~~~~~~~~~~~~~~~~~~~~~~~~~~~~ V frm
\begin{frame}[t,fragile] \frametitle{Pseudo-localization: the result}
  \lstterminal[basicstyle=\ttfamily\scriptsize]{codeoutput/Pseudo.txt}

  \vspace{0.1em}
  \small
  Hard-coded texts would stay plain (easy to spot); accents test fonts and encodings; \texttt{\textasciitilde\textasciitilde}
  makes texts longer to test the layout. \texttt{\{...\}} is kept, so values still work (texts inside \texttt{choice} stay plain).
\end{frame}
% ~~~~~~~~~~~~~~~~~~~~~~~~~~~~~~~~~~~~~~ A frm




% ~~~~~~~~~~~~~~~~~~~~~~~~~~~~~~~~~~~~~~ V frm
\begin{frame}[t,fragile] \frametitle{With Maven or Gradle}
  \begin{columns}[T]
    \begin{column}{0.48\textwidth}


% +++++++++++++++++++++++++++++++++++++++ V lst
%: -Lst.
\begin{lstlisting}[style=codesnippet, language={}, basicstyle=\ttfamily\scriptsize]
src/main/java/edu/yeditepe/shop/Main.java
src/main/resources/edu/yeditepe/shop/i18n/
    Messages.properties
    Messages_tr.properties
    Messages_de.properties
src/test/java/...      # tests in several locales
\end{lstlisting}
% +++++++++++++++++++++++++++++++++++++++ A lst


    \end{column}
    \begin{column}{0.48\textwidth}
      \begin{itemize}\small
        \item Resources are copied to the class path and packed into the JAR automatically.
        \item Set the encoding once:\\ \texttt{<project.build.sourceEncoding>UTF-8</...>}
        \item Run the tests with another locale:\\ \texttt{mvn test -DargLine="-Duser.language=tr"}
      \end{itemize}
      \small \hRemark{Not run here} (Maven needs the internet).
    \end{column}
  \end{columns}
\end{frame}
% ~~~~~~~~~~~~~~~~~~~~~~~~~~~~~~~~~~~~~~ A frm

% ~~~~~~~~~~~~~~~~~~~~~~~~~~~~~~~~~~~~~~ V
\hSection[Summary]{Summary}
% ~~~~~~~~~~~~~~~~~~~~~~~~~~~~~~~~~~~~~~ A




% ~~~~~~~~~~~~~~~~~~~~~~~~~~~~~~~~~~~~~~ V frm
\begin{frame}[t,fragile] \frametitle{Recipes}
  \begin{center}\scriptsize
  \begin{tabular}{L{0.3\textwidth}L{0.62\textwidth}}
    \toprule
    \textbf{goal} & \textbf{how} \\
    \midrule
    texts of a package     & \texttt{getBundle("edu.yeditepe.shop.i18n.Messages", l)} \\
    texts with values      & \texttt{new MessageFormat(rb.getString(key), l).format(values)} \\
    plurals                & \texttt{\{0,choice,0\#\dots|1\#\dots|1<\{0,number,integer\} \dots\}} \\
    money                  & \texttt{getCurrencyInstance(l)} + \texttt{setCurrency(...)}, \texttt{BigDecimal} \\
    dates                  & \texttt{ofLocalizedDate(FormatStyle.LONG).withLocale(l)}, then pass the \texttt{String} \\
    no default-locale fallback & \texttt{Control.getNoFallbackControl(Control.FORMAT\_DEFAULT)} \\
    safe parsing           & \texttt{parse(s, pos)} and check \texttt{pos.getIndex() == s.length()} \\
    search vs.\ sort       & \texttt{Collator} strength PRIMARY vs.\ TERTIARY \\
    bundle encoding        & UTF-8 (Java 9+), or \texttt{\textbackslash uXXXX} \\
    find missing texts     & pseudo-localization, tests in several locales \\
    \bottomrule
  \end{tabular}
  \end{center}
\end{frame}
% ~~~~~~~~~~~~~~~~~~~~~~~~~~~~~~~~~~~~~~ A frm




% ~~~~~~~~~~~~~~~~~~~~~~~~~~~~~~~~~~~~~~ V frm
\begin{frame}[t,fragile] \frametitle{Exercises}
  \begin{enumerate}\small
    \item Add \texttt{Messages\_fr.properties} and run \texttt{AllLocales} again.
    \item Give \texttt{Main} a \texttt{--currency EUR} option. What must change in the bundles? (Nothing!)
    \item Make the shop use \texttt{getNoFallbackControl} and repeat the trap experiment.
    \item Translate \emph{``\{0\} sent \{1\} files to \{2\}''} into Turkish and German with numbered placeholders.
    \item Write a test that fails when a key exists in \texttt{Messages} but not in \texttt{Messages\_tr}.
    \item Use \texttt{Collator} with PRIMARY strength to search names in a list without caring about case.
  \end{enumerate}
\end{frame}
% ~~~~~~~~~~~~~~~~~~~~~~~~~~~~~~~~~~~~~~ A frm




% ~~~~~~~~~~~~~~~~~~~~~~~~~~~~~~~~~~~~~~ V frm
\begin{frame}[t,fragile]
	\begin{center}
		\vspace{4cm}
		\hRemark{\Huge Questions?}
	\end{center}
\end{frame}
% ~~~~~~~~~~~~~~~~~~~~~~~~~~~~~~~~~~~~~~ A frm



\end{document}
